\chapter{Concept}
\label{conecpt}

In order to explore prompt engineering

\section{Automatic Prompt Optimization Using Gradient Descent}
Based on the work of \citewithauthor{pryzant2023AutomaticPrompt} I will implement a more sophisticated prompt optimization algorithm into the LiSSA framework. They propose the Prompt Optimization with Textual Gradients (ProTeGi) algorithm. This entire section is based on their work.

The ProTeGi algorithm takes an initial prompt $p$ and training data $\{(x_1, y_1), \dots, (x_n, y_n)\}$ consisting of input and output. They \directQuote[sec. 2]{assume access to a black box LLM API [...] which returns a likely text continuation y of the prompt formed by concatenating p and x}{pryzant2023AutomaticPrompt}. They then iteratively optimize the initial prompt p to produce an approximation of the most optimized prompt for the given task. In order to optimize the prompt, a function is required, to compute deviance between the actual output $y$ and expected output $y_i$ as a numeric value.

\begin{figure}[h]
\centering
\includegraphics[width=12cm]{logos/sdqlogo}
\caption{TODO: Overview of the iterative optimization loop}
\label{fig:sdqlogo}
\end{figure}

\textit{TODO: Continue explanation of APO with gradient descent}

